\documentclass[12pt]{article}
\usepackage{graphicx,amssymb,amsmath,amsthm}
\usepackage{hyperref,url}
\usepackage{verbatim}
\usepackage[labelfont=bf]{caption}
\newcommand{\nocopyright}{
No Copyright\thanks{
The authors hereby waive all copyright
and related or neighboring rights to this work,
and dedicate it to the public domain.
This applies worldwide.
}}
\title{Taut fillings}
\author{Peter Doyle \and Matthew Ellison \and Zili Wang\thanks{
Zili Wang is supported by
the National Natural Science Foundation of China (Grant No. 12301432)
}}
\date{Draft A}
\newtheorem{theorem}{Theorem}
\newtheorem{prop}{Proposition}
\newtheorem{lemma}[theorem]{Lemma}
\newtheorem*{hypo}{Hypothesis}
\newtheorem*{question}{Question}
\newtheorem{corollary}{Corollary}
\newtheorem*{conjecture}{Conjecture}

\def\startproof{{\bf {\medskip}{\noindent}Proof: }}
\def\definition{{\bf {\bigskip}{\noindent}Definition: }}
\def\endproof{$\spadesuit$  \newline}

\newcommand{\proofstart}{{\noindent \bf Proof.\ }}
\newcommand{\note}{{\noindent \bf Note.\ }}
\newcommand{\pfend}{\spadesuit}
\newcommand{\mathproofend}{\quad \qed}
\newcommand{\proofend}{$\quad \qed$}
\newcommand{\goesto}{\rightarrow}
\newcommand{\p}{\;\;\;}
\newcommand{\comma}{\p,}
\newcommand{\pd}{\p.}
\newcommand{\semi}{\p;}
\newcommand{\Rplus}{{\RR_{\geq 0}}}



\newcommand{\beginfigurepos}{\begin{figure}[htbp]}
%
% \fig
% followed by {filename}{label}{caption}
%
\newcommand{\fig}[3]{
\beginfigurepos
\includegraphics[width=370pt]{figures/#1.pdf}
\caption{#3}
\label{#2}
\end{figure}
}
%
% \figsize
% followed by {size}{filename}{label}{caption}
%
\newcommand{\figsize}[4]{
\beginfigurepos
\centerline{
\includegraphics[width=#1]{figures/#2.pdf}
}
\caption{#4}
\label{#3}
\end{figure}
}


\def\ZZ{\mathbf{Z}}
\def\QQ{\mathbf{Q}}
\def\RR{\mathbf{R}}
\def\heartsuit{\mbox{\boldmath{$Z$}}}% 
\def\del{\partial}
\def\dist{\mathrm{dist}}
\def\eps{\epsilon}


\newcommand{\parsec}{\subsection*}
\newcommand{\sect}[1]{\S\ #1}

%\newcommand{\gversus}[2]{\braket{#1\,|\,#2}}
\newcommand{\pair}[1]{\braket{#1}}
\newcommand{\sub}[1]{\braket{#1}}
\newcommand{\nnn}{\mathbf{NSet}}
\newcommand{\zzz}{\mathbf{ZSet}}
\newcommand{\nn}{\mathbf{N}}
\newcommand{\zz}{\mathbf{Z}}

\newcommand{\versus}[2]{\left(#1\,|\,#2\right)}
\newcommand{\gversus}[2]{\left[#1\,|\,#2\right]}
\newcommand{\poly}[1]{\left[#1\right]}

\newcommand{\rising}[1]{\overline{#1}}
\newcommand{\floor}[1]{\left \lfloor #1 \right \rfloor}
\newcommand{\ceiling}[1]{\left \lceil #1 \right \rceil}
\newcommand{\gbinom}{\genfrac{[}{]}{0pt}{}}

\makeatletter
\newcommand{\xRrightarrow}[2][]{\ext@arrow 0359\Rrightarrowfill@{#1}{#2}}
\newcommand{\Rrightarrowfill@}{\arrowfill@\equiv\equiv\Rrightarrow}
\newcommand{\xLleftarrow}[2][]{\ext@arrow 3095\Lleftarrowfill@{#1}{#2}}
\newcommand{\Lleftarrowfill@}{\arrowfill@\Lleftarrow\equiv\equiv}
\newcommand{\xLleftRrightarrow}[2][]{\ext@arrow 3399\LleftRrightarrowfill@{#1}{#2}}
\newcommand{\LleftRrightarrowfill@}{\arrowfill@\Lleftarrow\equiv\Rrightarrow}
\makeatother

\newcommand{\eq}{\Rrightarrow}
\newcommand{\funeq}[1]{\stackrel{#1}\eq}
\newcommand{\funequiv}[1]{\stackrel{#1}\equiv}
\newcommand{\funto}[1]{\stackrel{#1}\to}
\newcommand{\union}{\cup}
\newcommand{\cross}{\times}
\newcommand{\op}{\operatorname}
\newcommand{\kw}{\mathrm}
\newcommand{\escape}{\op{escape}}
\newcommand{\cancel}{\op{cancel}}
\newcommand{\nestuntil}{\operatorname{nestuntil}}
\newcommand{\size}{\operatorname{size}}
\newcommand{\id}{\operatorname{id}}
\newcommand{\create}{\operatorname{create}}
\newcommand{\destroy}{\operatorname{destroy}}
\newcommand{\then}{\triangleleft}
\newcommand{\lp}{\kw{LP}}
\newcommand{\LPvol}{\kw{LPvol}}
\newcommand{\vol}{\kw{vol}}
\newcommand{\flipdist}{\kw{flipdist}}
\newcommand{\tetbar}{\kw{homvol}}
\newcommand{\vsa}{\kw{vsa}}
\newcommand{\cone}{\kw{cone}}
\newcommand{\link}{\kw{link}}
\newcommand{\maxdeg}{\kw{maxdeg}}
\newcommand{\mindeg}{\kw{mindeg}}
\newcommand{\vertices}{\kw{Vert}}
\newcommand{\carrier}{\kw{carrier}}
\newcommand{\nbhd}{\kw{nbhd}}
\newcommand{\adj}{\kw{adj}}
\newcommand{\taut}{\kw{taut}}
\newcommand{\net}{\kw{net}}

\newcommand{\Zvol}{\kw{Zvol}}
\newcommand{\Qvol}{\kw{Qvol}}
\newcommand{\Rvol}{\kw{Rvol}}
\newcommand{\tetvol}{\kw{tetvol}}
\newcommand{\hash}{\#}
\newcommand{\ms}{[}
\newcommand{\msend}{]}

\newcommand{\kwstyle}{\bf}
\newcommand{\IF}{{\ \kwstyle{if}\ }}
\newcommand{\ELSE}{{\ \kwstyle{else}\ }}
\newcommand{\WHILE}{{\kwstyle{while}\ }}
\newcommand{\DO}{{\ \kwstyle{do}\ }}
\newcommand{\RETURN}{{\kwstyle{return}\ }}
\newcommand{\direct}{\oplus}
\newcommand{\pls}{\mathord{+}}
\newcommand{\mns}{\mathord{-}}



\begin{document}

\maketitle

\begin{abstract}
Let $\sigma$ be a simplicial triangulation of the 2-sphere,
$X$ the associated integral 2-cycle.
A filling of $X$ is an integral 3-chain $M$
with $\del M = X$;
a taut filling is one with minimal $L_1$-norm.
We show that any taut filling 
arises from an extension of $\sigma$ to
a simplicial complex homeomorphic to the 3-ball.
This complex is shellable, and `flag':
every clique in its 1-skeleton occurs as a simplex.
The key to the proof is the general
fact that any taut filling of an $n$-cycle
splits under disjoint union, connected sum, and more generally
what we call almost disjoint union,
where summands are supported on sets that overlap in at most $n+1$ vertices.
\begin{comment}
Despite the generality of this result, we have nothing to say about
optimal fillings of spheres of dimension 3 or higher.
\end{comment}
\end{abstract}

\section{Overview}
Let $\Delta = \Delta(\Omega)$ be the $|\Omega|\mns 1$-simplex
with vertices $\Omega$,
viewed as the abstract simplicial complex consisting of all subsets of
$\Omega$.
Let $C_n=C_n(\Omega)=C_n(\Delta(\Omega),\ZZ)$
and $Z_n=Z_n(\Omega)=Z_n(\Delta(\Omega),\ZZ)$
be its integral $n$-chains and $n$-cycles.
A \emph{filling} of an $n$-cycle $X \in Z_n$ is any
$n\pls 1$-chain $M \in C_{n+1}$ with $\del M = X$.
(Here and throughout we'll assume $n\geq 1$.)
Let $\Zvol(X)$ be the minimum $L_1$-norm of a filling of $X$.
Call $M \in C_{n+1}$ \emph{taut} if it is an optimal filling
of its boundary, i.e., if $|M|=\Zvol(\del M)$,

For $X \in C_n$ let
$\vertices(X)$ be the set of all the vertices of all the $n$-simplices
to which $X$ assigns non-$0$ weight.
For any taut filling $M$ of $X$ we have $\vertices(M) = \vertices(X)$.
(Cf.\ Proposition \ref{nointernal} below.)
This is why we write $C_n$ and $Z_n$ without reference to $\Omega$.

Call an $n$-cycle $X+Y \in Z_n$ an \emph{almost disjoint union} if
\[
|\vertices(X) \cap \vertices(Y)| \leq n+1
.
\]
This notion generalizes both disjoint union and connected sum
along an $n$-simplex.
Ellison \cite{ellison:gap}
showed the $\Zvol$ adds under almost disjoint union:
$\Zvol(X+Y) = \Zvol(X) + \Zvol(Y)$.
This means that \emph{some} taut filling $M$ of $X$ splits into a sum
$M=M_X+M_Y$ of taut fillings of $X,Y$.
Here we amplify Ellison's Corollary 2 to show (Theorem \ref{th1})
that when $n \geq 2$, \emph{any} taut filling of $X+Y$ splits.

In place of integral chains we can use chains with
coefficients in $\QQ$ or $\RR$.
Evidently $\Rvol \leq \Qvol$.
But computing $\Rvol$ is
a linear programming problem with rational coefficients,
so in fact $\Qvol = \Rvol$.
Ellison used LP duality to prove that
$\Qvol$ adds under almost
disjoint union.
Now for $n \geq 2$ we get the stronger result
(Corollary \ref{cor1})
that taut $\QQ$-fillings split:
Multiplying a taut
$\QQ$-filling $M$ of $X$ by a common denominator $q$
yields a taut $\ZZ$-filling $qM$ of $qX$:
Split the $\ZZ$-filling $qM$ and divide by $q$ to get a splitting
of the $\QQ$-filling $M$.

Now let $\sigma$ be a simplicial triangulation of $S^2$;
let $X(\sigma) \in Z_2(\Delta(\vertices(\sigma)),\ZZ)$
be either one of the 2-cycles
that arises from $\sigma$ by orienting its $2$-simplices;
and write $\Zvol(\sigma) = \Zvol(X(\sigma))$.
Let $\tetvol(\sigma)$ be the minimum number of $3$-simplices
required to extend $\sigma$ to a simplicial triangulation $\tau$ of $B^3$.
We will show that
\[
\Zvol(\sigma) = \tetvol(\sigma)
.
\]
Certainly $\Zvol \leq \tetvol$,
because any extension $\tau$ induces a filling
of $X(\sigma)$.
Theorem \ref{th2} states that 
any taut filling of $X(\sigma)$ arises in this way.
The proof of Theorem \ref{th2} relies on Theorem \ref{th1}.
\begin{comment}
We use induction, with base case $\sigma$ a tetrahedron.
In any taut filling $Y$ of $X(\sigma)$,
thought of as a multiset of oriented $3$-simplices,
some $t \in Y$ meets $\sigma$ in at least two faces.
Ignoring the base case,
$Y-t$ is a taut filling of
$X(\sigma) - \del t = X(\sigma^\prime)$,
where $\sigma^\prime$ is either a simplicial triangulation of $S^2$,
or two such triangulations $\sigma_1,\sigma_2$ fused along an edge.
In the first case induction yields a triangulation of $B^3$,
and we glue on $t$.
In the second case by Theorem \ref{th1}
$Y-t$ splits into taut fillings of $X(\sigma_1),X(\sigma_2)$;
by induction both give triangulations of $B^3$, which we
glue onto $t$.
So we wind up with some kind of triangulation $\tau$ of $B^3$,
built from the simplices of $Y$.
The wrinkle is that we need to ensure that in $Y$ there are no
identifications beyond those in $\tau$.
\end{comment}

We don't know what happens for spheres of dimension 3 or greater.

\section{Background}

Our interest in $\Zvol$ stems from the work
of Sleator, Tarjan, and Thurston
\cite{stt:jams}.
They observed that coning from a vertex of maximal degree
shows that a $v$-vertex
triangulation $\sigma$
of $S^2$ has $\tetvol \leq 2v-4-\maxdeg(\sigma)$.
If $v\geq13$ $\maxdeg(\sigma) \geq 6$, so $\tetvol \leq 2v-10$.
They
produced examples for which $\tetvol = 2v-10$, provided
$v$ is larger than some unspecified bound,
and conjectured that such examples should exist for any $v \geq 13$.
Their approach was to 
realize $\sigma$ as an ideal hyperbolic polyhedron
with volume $2v V_0-O(\log(v))$,
where $V_0$ is the volume of an
equilateral ideal tetrahedron,
which is maximal.
This implies $\tetvol \geq 2v-O(\log(v))$.
From there they bootstrapped their way up to showing
that $\tetvol=2v-10$.
They suggested
\cite[p. 697]{stt:jams}
that in fact $\Qvol= 2v-10$,
which would immediately imply $\tetvol=2v-10$.
Mathieu and Thurston \cite{mt}
produced a class of examples,
different from the examples of STT,
for which they could show that $\Qvol = 2v-10$,
still under the assumption that $v$
is sufficiently large.
In \cite{dew:filling} we produced examples for any $v \geq 13$ with
$\Qvol=2v-10$, confirming the conjecture of Sleator, Tarjan, and Thurston.
Theorem \ref{th2} here
tells us that whenever we have $\Qvol=2v-10$,
or just $\Zvol=2v-10$, any taut filling
must arise from a triangulation of the ball.

About $\Qvol$ versus $\Zvol$.
In \cite{dew:filling}
we described triangulations of $S^2$ with $\Qvol < \Zvol$.
In all such examples that we know of
with $\maxdeg \leq 6$, the gap between $\Qvol$ and $\Zvol$ is $<1$,
so that
$\lceil \Qvol \rceil = \Zvol$.
Since $\Qvol$ and $\Zvol$
add under connected sum,
the gap between them can be arbitrarily
large
(Ellison \cite{ellison:gap}),
but taking the connected sum pushes $\maxdeg$ above $6$.



\section{Taut fillings}

We can think of an $n$-chain $X \in C_n$ as a multiset of
non-cancelling oriented simplices:
\[
X = \sum_{t \in X} t
.
\]
In this sum $t$ denotes an oriented simplex
\[
t=[x_0,\ldots,x_n]
= [x_{\pi(0)},\ldots,x_{\pi(n)}]
,\;\;\mbox{$\pi$ an even permutation}
,
\]
and each unoriented simplex contributes a number of terms corresponding
to its multiplicity.

The size of $M$ as a multiset is its $L_1$-norm $|M|$.
Write $U \subset M$
if $U$ is a sub-multiset of $M$.
This happens just when
$|M| = |U| + |M-U|$.

\begin{prop} \label{subtaut}
If $M$ is taut and $U \subset M$ then $U$ is taut.
\end{prop}

\proofstart
If $U$ is not taut, there is a smaller filling $V$ of $\del U$.
Take $M^\prime = M-U+V$.
In terms of multisets, this means that we substitute $V$ for $U$,
and then do any required cancellation of oppositely oriented simplices
of $M-U$ and $V$.
We have
\[
\del M^\prime = \del M - \del U + \del V = \del M
\]
and
\[
|M^\prime| \leq |M-U| + |V| = |M| - |U| + |V| < |M|
,
\]
contradiction.
\proofend

\section{Coning}
For $x \in \Omega$, $U \in C_n$ let
\[
\nbhd(x,U) = \sum_{ t \in U: x \in \vertices(t) } t
,
\]
\[
\deg(x,U) = |\nbhd(x,U)|
,
\]
\[
\maxdeg(U) = \max_x \deg(x,U)
.
\]
The \emph{cone from $x$ to $U$}
is the $(n+1)$-chain
\[
\cone(x,U) = 
\sum_{t \in U} \adj(x,t)
,
\]
consisting of all the non-trivial oriented
$(n+1)$-simplices $\adj(x,t)=[x x_0\ldots x_n]$ obtained
by adjoining $x$ to  $t = [x_0\ldots x_n] \in U$.
If $x \in \vertices(t)$ then $t$ doesn't contribute to the sum,
so $|\cone(x,U)| = |U| - \deg(x,U)$.

If $U$ is closed then
$\del\, \cone(x,U) = U$, so

\begin{prop} \label{maxdeg}
For any $X \in Z_n$
\[
\Zvol(X) \leq |X| - \maxdeg(X)
\mathproofend
\]
\end{prop}

If $U \in Z_n$ and $x \notin \vertices(U)$ then
$\deg(x,U)=0$ and
$|\cone(x,U)| = |U|$.
In this case we call $\cone(x,U)$ a \emph{complete cone}.
By Proposition \ref{maxdeg} a non-trivial complete cone is not taut.
(Equivalently: rather than coning from $x \notin U$, one could cone from
some $x \in U$.)
Thus:
\begin{prop}\label{nocone}
If $M$ is taut it contains no non-trivial complete cone.
\proofend
\end{prop}

Call $x \in \vertices(M) \setminus \vertices(\del M)$ an
\emph{internal vertex} of $M$.
If $x$ is internal to $M$ then $\nbhd(x,M)$ is a complete cone, so:
\begin{prop} \label{nointernal}
If $M$ is taut then
it has no internal vertices.
\proofend
\end{prop}

\section{Almost disjoint unions}

Recall that if $X,Y \in Z_n$ we call $X+Y$ an almost disjoint union
if
\[
|\vertices(X) \cap \vertices(Y)| \leq n+1
.\]
The most interesting special case is a connected sum,
where $t=[c_0,c_1,\ldots,c_n]$ occurs once in $X$ and
$-t=[c_1,c_0,\ldots,c_n]$ occurs once in $Y$.
For example, if $n=1$ with $X$ a cycle of length $p$
and $Y$ a cycle of length $q$
the connected sum $X+Y$ is
a cycle of length $p+q-2$.
Here
$\Zvol$ adds, because
$\Zvol(X)=p-2$,
$\Zvol(Y) = q-2$,
and
\[
\Zvol(X+Y) = (p+q-2)-2 = \Zvol(X)+\Zvol(Y)
.
\]
But not every filling of $X+Y$ splits, because a taut filling
may contain any 2-simplex with vertices in $\vertices(X+Y)$.

We want to show that fillings do split when $n \geq 2$.

For $A \subset \Omega$, $p \in A$ define
\[
\pi_{A,p}: \Omega \to A
\]
\[
\pi_{A,p}(x) = x \IF x \in A \ELSE p
\]
and let
\[
K_*(A,p): C_*(\Omega) \to C_*(A)
\]
be the induced chain map.
This is a projection of $C_*(\Omega)$ onto $C_*(A)$.

Let $A,B$ be finite vertex sets sharing the vertices
$C = A \cap B$,
and let $c=|C|$.
Observe that
$C_c(C)$ and $Z_{c-1}(C)$ are trivial,
because $\Delta(C)$ contains no $c$-simplex,
and (up to orientation) only one $c\mns 1$-simplex,
with nothing to cancel its boundary.

Let
\[
C_*(A,B) = C_*(A) \direct C_*(B)
.
\]
For $p,q \in C$ define the mapping
\[
g_*(p,q) = K_*(A,p) \direct K_*(B,q):
C_*(A \cup B) \to C_*(A,B)
.
\]

\begin{prop} \label{recover}
Let $(X,Y) \in C_n(A,B)$.
If $|A \cap B| \leq n$
we can recover $(X,Y)$ from $X+Y$.
If $(X,Y) \in Z_n(A,B)$ this holds also when $|A \cap B|  = n+1$.
\end{prop}

\proofstart
Let $C=A \cap B$.
We can assume $|C| \geq 1$.
(Add a brand new point to $C$ if necessary.)
We claim that
for any $p,q \in C$
(not necessarily distinct) we have
\[
g_n(p,q)(X+Y)=(X,Y)
.
\]
Indeed,
\[
K_n(A,p)(X+Y) = X + K_n(A,p)(Y)
.
\]
The second term belongs to $C_n(C)$,
which is trivial when $|C| \leq n$.
If $Y \in Z_n(B)$ the second term belongs to $Z_n(C)$,
which is trivial when $|C| \leq n+1$.
\proofend


\begin{theorem} \label{th1}
If $|A \cap B| \leq n+1$,
for all $(X,Y) \in Z_n(A,B)$ we have
\[
\Zvol(X+Y) = \Zvol(X)+\Zvol(Y)
.
\]
And as long as $n \geq 2$,
for any $M \in \taut(X+Y)$ 
we have $M = M_X + M_Y$
with $M_X \in \taut(X)$, $M_Y \in \taut(Y)$.
\end{theorem}

\note
Ellison \cite{ellison:gap}
proved the additivity of $\Zvol$.
What's new here is the splitting of taut fillings when $n\geq 2$.
This result resembles Theorem 6.2 of
Pournin and Wang
\cite{pw:flip}
about flip paths.
Like theirs, our proof uses a variation on the normalization technique
of STT
\cite[Lemma 7]{stt:jams}.
It would be nice to fit these results under one roof.
\medskip

\proofstart
Again let $C = A \cap B$.
We can assume $|C| = n+1$,
as this is the hardest case.
And we might as well go ahead and take $n=2$, $|C|=3$,
as this case illustrates all the issues.

Take any $M \in \taut(X+Y) \subset C_3(A \union B)$.
Pick distinct points $p,q \in C$,
and let
\[
(M_X,M_Y) = g_{3}(p,q)(M)
.
\]
(For now we'll suppress the dependence of $M_X,M_Y$ on $p,q$.)

Because $g_*(p,q)$ is a chain map we have
$\del_{3} M_X = X$, $\del_{3} M_Y = Y$:
\[
(\del_{3} M_X, \del_{3} M_Y)
=
\del_{3} (g_{3}(p,q)(M))
=
g_2(p, q)(\del_{3} M)
=
g_2(p,q)(X+Y)
=
(X,Y)
.
\]

We want to show that under the map $g_3(p,q)$ any $t \in M$
dies either in $M_X$ or $M_Y$,
because then we'll get additivity of $\Zvol$ from
\[
\Zvol(X+Y) = |M|
\geq |M_X| + |M_Y|
\geq \Zvol(X) + \Zvol(Y)
\geq \Zvol(X+Y)
.
\]

Let's call a $3$-simplex a `tet' for short.
Say that a tet $t \in M$ has type $CCXY$ if $t= \pm[c_1c_2x_1y_1]$
for $c_1,c_2 \in C$, $x_1 \in A \setminus C$, $y_1 \in B \setminus C$.
Similarly for types $XXXX$, $CXXX$, $CXXY$, $XXXY$, etc.
The first two are pure $X$ cases, meaning that they live in $C_3(A)$;
the last two are hybrid cases.

Any $XX..$ tet dies in $M_Y$;
any $YY..$ tet dies in $M_X$.

The remaining cases to check are
are the pure cases $CCCX$, $CCCY$,
and the hybrid case $CCXY$.
A $CCCX$ tet $t=\pm[c_1c_2c_3x_1]$ dies in $M_Y$ because
$K_3(B,q)(t)=[c_1c_2c_3q]$,
which vanishes because $q \in C = \{c_1,c_2,c_3\}$ produces a repeated vertex.
Likewise a $CCCY$ tet dies in $M_X$.
The more interesting case is $CCXY$:
The key is that since $|C|=3$, $\{c_1,c_2\}$ cannot be disjoint from $\{p,q\}$,
so it dies on one side or the other.

So 
\[
|M| = |M_X(p,q)| + |M_Y(p,q)|,
\]
and $\Zvol$ adds.
Note how we're now emphasizing the possible dependence of $M_X,M_Y$ on $p,q$.
When $n=1$ the choice of $p,q$ can indeed make a difference:
We can get a different pair $(M_X,M_Y)$ if we switch
$p$ and $q$.
(Think about the connected sum of two cycle graphs.)

But when $n \geq 2$ we will now show that all tets are pure $X$ or pure $Y$,
so $M_X$ consists of all the pure $X$ tets, and ditto for $M_Y$.
This will make
$M = M_X + M_Y$ with $M_X \in \taut(X)$, $M_Y \in \taut(Y)$.

Again our test case is $n=2$.

The key observation is that, now that we know that every tet dies on
one side or the other, we know that no tet can die on both sides,
because that would make $|M| > \Zvol(X) + \Zvol(Y)$.

An $XXYY$ tet would die on both sides, so there can be none of these.

Any $pqXY$, $pXXY$, or $qXYY$ would die on both sides, and $p,q$ are arbitrary,
so this rules out all $CCXY,CXXY,CXYY$.

The only remaining hybrids are $XXXY$ and $XYYY$.
Let's rule out $XXXY$, leaving $XYYY$ to symmetry.

Fix any $y_0 \in B \setminus C$, and suppose there is
some $XXXy_0$ tet in $M$.
Let $U \subset M$ consist of all such $XXXy_0$ tets.
We claim that $y_0 \notin \vertices(\del U)$.
For let $s$ be any 2-simplex of the form $XXy_0$,
the only kind of 2-simplex containing $y_0$ that could belong
to $\del U$.
Any $t \in M$ of which $s$ or $-s$ is a face must be a hybrid tet,
and the only possibility is $XXXy_0$, so $t \in U$.
Because $\del M = 0$ the signed multiplicity of $s$
vanishes in $\del M$,
hence also in $\del U$,
because the tets of $M$ with $\pm s$ as a face all belong to $U$.
So $y_0 \notin \vertices(\del U)$,
making $U$ a complete cone on $y_0$,
contradicting \ref{nocone}.
So there is no such $XXXy_0$ tet, ruling out $XXXY$,
and with it $XYYY$.

So there are no hybrid tets, meaning that $M$ splits, as claimed.
\proofend

The splitting of taut fillings depends on the assumption $n\geq 2$.
Let's consider what goes wrong when $n=1$.
Here the hybrid types are $CXY$, $XXY$, $XYY$.
Neither $pXY$ nor $qXY$ dies on both sides, so we can't rule out $CXY$.
Failing that, the $XXy_0$ tets needn't form a complete cone,
because $[x_0x_1y_0)$ can continue across $[x_1 y_0]$ to $-[p x_1 y_0]$,
so we can't rule out $XXY$ either.

The foregoing proof works equally well over $\QQ$, providing we
permit ourselves to work with fractional multisets.
Alternatively, we can clear denominators as in the introduction
above.
Either way, we have:

\begin{corollary} \label{cor1}
$\Qvol$ adds under almost disjoint union,
and for $n \geq 2$
taut $\QQ$-fillings split.
\proofend
\end{corollary}

\section{Triangulations}

We turn now to filling simplicial triangulations of $S^2$.
Let's begin by establishing terminology.

An \emph{$n$-simplex} $s$ is simply a set of size $n+1$.
Its \emph{faces} are its subsets,
which are $k$-simplices with $-1 \leq k \leq n$,
$-1$ being the dimension of the empty simplex.
A \emph{simplicial complex} $\sigma$ is a finite collection of simplices
closed under taking faces:
If $t \subset s \in \sigma$ then $t \in \sigma$.

For any $k$-simplex $s \in \sigma$ define the \emph{link}
\[
\link(s,\sigma) = \{t: t \cap s = \emptyset, s \cup t \in \sigma \}
.
\]
This is a simplicial complex, but
(except for $\link(\emptyset,\sigma)=\sigma$)
it is not a subcomplex of $\sigma$,
because we are taking simplices in $\sigma$ and knocking down
their dimension by $k+1$.

We are interested in particularly nice simplicial complexes
called normal pseudomanifolds, or as we will prefer to say,
\emph{clean $n$-complexes}.

\begin{enumerate}
\item
To start with,
a clean complex must be \emph{pure},
meaning that every simplex of $\sigma$ belongs to some $n$-simplex in $\sigma$.
To determine $\sigma$ we can specify its $n$-simplices,
and then throw in all their subsets.
So we can confound a pure complex with the set of its $n$-simplices.

\item
A clean complex is a \emph{pseudomanifold}:
Every $n$-simplex abuts at most one other $n$-simplex
across any given $n-1$-simplex.
Another way to say this is that
the link $\link(s,\sigma)$ of any $n-1$-simplex $s$
consists of either a single point 
(in which case $s \in \del \sigma$ is a \emph{boundary simplex}),
or two points
($s \in \sigma \setminus \del \sigma$ is an \emph{interior simplex}.)

\item
A clean complex is \emph{normal}:
For any simplex $s$ of dimension $0,\ldots,n-2$,
$\link(s,\sigma)$ is connected.
(Some definitions extend this requirement
to the empty simplex of dimension $-1$;
this forces $\sigma$ to be connected.)
Being normal means that $\sigma$ is just what you get by gluing its 
$n$-simplices along shared faces,
without extra identifications. This rules out, for example,
an icosahedron with a pair of opposite vertices identified.
\end{enumerate}

If $\sigma$ is clean, 
$\link(s,\sigma)$ is clean.
The boundary $\del \sigma$ of $\sigma$ is clean,
and closed:
$\del \del \sigma = \{\emptyset\}$.

If $W$ is an $n$-manifold,
say that a simplicial complex $\sigma$
is a \emph{simplicial triangulation} of $W$
if the geometric carrier of $\sigma$ is homeomorphic to $W$.
In this case $\sigma$ is necessarily a clean complex.
We imagine that we have prescribed a specific homeomorphism,
at least up to the point of picking out one particular orientation
$X(\sigma) \in C_n$ when $W$ is oriented;
but we do not insist upon this, because nothing will depend on which
orientation we pick.

This notion of triangulation is not as general as you might want
for some purposes.
The double cover of a triangle is not a simplicial triangulation
of $S^2$,
because that would requires two distinct
$2$-simplices to have the same three edges.
Nor can any $2$-simplex have two of its edges
glued to one another.
Thurston
\cite{thurston:shapes}
allows such triangulations,
but it is unclear how important these are to
his theory of shapes of surfaces.
We don't allow them.

We continue to think of a chain $M \in C_n$ as a multiset of
oriented $n$-simplices.
We call $M$ \emph{simplicial} if no simplex is repeated, that is, if every
simplex occurs in $M$ with multiplicity $0$ or $\pm 1$.
In that case we write $K(M)$ for the \emph{support} of $M$: the pure
simplicial complex generated by the $n$-simplices occurring in $M$.
(In passing to $K(M)$ we view oriented simplices as a subclass of
simplices, and identify a set of $n$-simplices with the pure $n$-complex
it generates.)
We call $M$ \emph{clean} if it is simplicial and $K(M)$ is clean.

In particular, a filling $M$---an $(n\pls1)$-chain---is \emph{clean} when
it has no repeated simplices and its support $K(M)$ is a clean simplicial
complex. For a filling of a triangulated $2$-sphere, the simplices in
question are tetrahedra.

\section{Filling a triangulation of the 2-sphere}

\begin{theorem} \label{th2}
Let $\sigma$ be a simplicial triangulation of $S^2$,
and $M$ a taut filling of $\sigma$.
Then $M$ is clean, and its support $K(M)$
is a simplicial triangulation of $B^3$.
\end{theorem}

\proofstart

Let $v,e,f$  count the vertices, edges, and faces of $\sigma$.
We have
$e=3v-6$, $f=2v-4$.
(Check: $3f=2e$, $v-e+f=2$.)

Consider a counterexample pair $(\sigma,M)$ for which
$\Zvol(\sigma)=|M|$ is minimal.
Obviously $v>4$.
We can assume that $\sigma$ is prime
(not a connected sum along a triangle),
and in particular (this is all we will need) that there is no
vertex of degree $3$.

Call a tet $t \in M$ \emph{eligible} if it shares two of its oriented
faces with $\sigma$. (It cannot share more than two, since $\sigma$ has
no vertex of degree $3$.)
By Proposition \ref{maxdeg},
\[
|M| = \Zvol(\sigma) \leq f-\maxdeg(\sigma)
.
\]
Regard $M$ as a multiset of tets, and assign to each $2$-simplex
$\alpha$ of $\sigma$ some tet of $M$ having $\alpha$ as a properly
oriented face. Such a tet exists, because $\alpha$ occurs with
coefficient $+1$ in $\del M = X(\sigma)$, so some tet of $M$ contributes
$\alpha$ with the correct sign. (Throughout we count support tets and
correct boundary signs, not coefficient copies or cancelled appearances.)
This assignment is at most $2$-to-$1$, since no tet of $M$ has more than
two faces on $\sigma$, so at least $\maxdeg(\sigma)$ tets receive two
faces each. These are \emph{disjointly eligible}: eligible tets whose
pairs of boundary faces are mutually disjoint. We will need only two.

Consider an eligible tet $t$ with boundary faces $s_1=[a,b,c]$ and
$s_2=[b,a,d]$. Removing $t$ from $M$ produces a filling $M-t$ of a
$2$-complex $\sigma_t$, in which $s_1,s_2$ have been replaced by the
faces $[c,d,b],[d,c,a]$: the edge $ab$ has been \emph{flipped} to the
edge $cd$. There are two cases.
\begin{itemize}
\item[(1)] Edge $cd$ was not already an edge of $\sigma$. Then
$\sigma_t$ is again a simplicial triangulation of $S^2$, and $M-t$ is a
taut filling of it. By minimality of $|M|$, $M-t$ is clean.
\item[(2)] Edge $cd$ was already an edge of $\sigma$. Then $\sigma_t$
is the almost disjoint union of two triangulations $\sigma_1,\sigma_2$
of $S^2$, conjoined along the edge $cd$. By Theorem \ref{th1}, $M-t$
splits as $M-t=M_1+M_2$, where each $M_i$ is a taut filling of
$\sigma_i$; by minimality, each $M_i$ is clean. Here $M-t$ itself need
\emph{not} be clean as a single complex---the link of the edge $cd$ in
$K(M-t)$ is disconnected---even though it splits into the clean pieces
$M_1,M_2$. It is, however, still simplicial.
\end{itemize}
In either case $M-t$ is simplicial.

We must now show that $M$ itself is clean. First, $M$ is simplicial: the
only way it could fail to be is for some eligible tet $t$ to occur in $M$
with multiplicity $2$. Choose a second tet $u$, disjointly eligible from
$t$, so that $u\neq t$. Removing $u$ leaves $2t$ present in $M-u$, so
$M-u$ is not simplicial---contradicting the previous paragraph applied to
$u$. Hence $M$ is simplicial.

It remains to verify that $M$ is clean, by checking the three conditions
defining a clean complex.

\begin{enumerate}

\item
$M$ is a pseudomanifold.
For suppose some $2$-simplex $s$ occurs more than twice as the boundary
of a $3$-simplex of $M$.
If $s \notin \sigma$,
it must occur at least twice plus and twice minus;
after removing any eligible tet it
still occurs either twice plus or twice minus in $M-t$ in case (1),
or in one or the other of $M_1,M_2$ in case (2),
contradicting minimality.
If $s \in \sigma$,
it occurs at least twice plus in $M$:
Remove some eligible tet that does not have $s$ as a face
to get a contradiction.

\item
$M$ has no edge $\{a,b\}$ whose link
is disconnected.
For the link is a 1-dimensional complex whose edges correspond to
tets of $M$ that contain $\{a,b\}$.
The only way a component of the link can have fewer than three
vertices is if consists of a single edge $\{c,d\}$, corresponding
to a tet $t=\{a,b,c,d\} \in M$.
In this case $\{a,b\}$ must be an edge of $\sigma$,
and $\{a,b,c\},\{a,b,d\}$ the faces of $\sigma$
that are adjacent along $\{a,b\}$.
These are the only faces of $\sigma$ that contain $\{a,b\}$,
so no other component of the link can reduce to a single edge.
If the link is disconnected, removing any eligible tet
other than $t$ will
leave it disconnected, contradicting minimality.

\item
$M$ has no vertex whose link is disconnected:
For if $\link(\{a\},M)$ is not connected,
the only way removing a single tet $t$
can render the link connected
is if one component of the link has a single $2$-simplex
$\{b,c,d\}$ and $t=\{a,b,c,d\}$.
In this case $\{a,b,c\}$, $\{a,b,d\}$, $\{a,c,d\}$ must
all be faces of $\sigma$,
making $a$ a vertex of $\sigma$ of degree $3$, contradiction.

\item Hence $M$ is clean.
\proofend

\end{enumerate}

Thanks to this theorem we may, when convenient, speak of a taut filling
$M$ of $\sigma$ as `being' the simplicial triangulation $\tau=K(M)$ of
$B^3$.

\section{Shelling}

A \emph{shelling} of a simplicial triangulation $\tau$ of $B^3$ is an
ordering $(t_1,\ldots,t_{|\tau|})$ of its tets such that gluing them on
one at a time maintains a triangulation of $B^3$: for each
$k=1,\ldots,|\tau|$, the subcomplex $\tau_k$ generated by
$\{t_1,\ldots,t_k\}$ is a simplicial triangulation of $B^3$.
For $k>1$ we classify $t_k$ by the number of faces it shares with
$\tau_{k-1}$. Gluing along one face increases the vertex count $v$ by
$1$; along two faces leaves $v$ unchanged; along three faces decreases
$v$ by $1$. In our setting the third type never occurs, since gluing a
tet along three faces would bury an interior vertex; thus $v$ never
decreases. (A shelling with this property is sometimes called
\emph{monotone}; here every shelling is monotone, so we do not belabor
the distinction.)

Call $\tau$ \emph{freely shellable} if any $t \in \tau$ can serve as the
initial tet of a shelling.

\begin{theorem} \label{th3}
If $\sigma$ is a simplicial triangulation of $S^2$, then any taut filling
$M$ of $\sigma$ is freely shellable: its support $\tau=K(M)$ is a
simplicial triangulation of $B^3$ admitting a shelling that begins with
any prescribed tet.
\end{theorem}

\proofstart
Fix the designated initial tet $t_1 \in \tau$. Rather than building
$\tau$ up starting from $t_1$, we describe a recursive procedure---call
it \emph{shucking}---that takes $\tau$ apart so as to remove $t_1$ last;
reversing it yields a shelling that begins with $t_1$.

The base case is $\tau = \{t_1\}$.

If $\tau$ has a vertex of degree $3$ lying on a tet $t \neq t_1$, remove
$t$ and shuck $\tau-t$, ending with $t_1$.

If $\tau$ has a vertex of degree $3$ lying on $t_1$, let $t_2$ be the tet
of $\tau-t_1$ to which $t_1$ is attached; shuck $\tau-t_1$ ending with
$t_2$, then remove $t_1$.

If $\tau$ has no vertex of degree $3$, then, as in the proof of
Theorem \ref{th2}, there are at least $\maxdeg \geq 4$ disjointly eligible
tets. Pick one, $t$, sharing no face with $t_1$. If $\tau-t$ does not fall
apart, remove $t$ and shuck $\tau-t$ ending with $t_1$. If $\tau-t$ falls
apart into the almost disjoint union of $\tau_1$ and $\tau_2$, with
$t_1 \in \tau_1$ and $t$ bridging the two pieces, let $t_2$ be the tet of
$\tau_2$ to which $t$ attaches; shuck $\tau_2$ ending with $t_2$, remove
$t$, then shuck $\tau_1$ ending with $t_1$.
\proofend

\section{Flag complex}

A simplicial complex is a \emph{flag complex} (also called a
\emph{clique complex}) if every clique in its $1$-skeleton occurs as a
simplex.
A simplicial triangulation $\sigma$ of $S^2$
is flag just if it is \emph{prime} 
(not the connected sum of smaller complexes)
and not the boundary of a tetrahedron.

We can decompose any $\sigma$ as the connected sum of
prime components $\sigma_i$.
(Of course, to reconstruct $\sigma$ we need more than this, we need to know
just how they are connected.)
From the point of taut fillings these components are independent:
Any taut filling $\tau$ of $\sigma$ is a union of taut fillings $\tau_i$.
That each of these is freely shellable implies the same of the union.

Now we want to show that, whether or not $\sigma$ itself is prime, 
any taut filling of $\sigma$ is a flag complex.

\begin{theorem}
If $\sigma$ is a simplicial triangulation of $S^2$,
any taut filling $\tau$ of $\sigma$ is a flag complex.
\end{theorem}

\proofstart
The goal here is to show that neither of these
three taboo configurations occurs
in $\tau$:
\begin{enumerate}
\item
The three edges of triangle $\{A,B,C\}$ without the 2-simplex $ABC$.
\item
The six edges of the $K_4$ $\{A,B,C,D\}$ without the 3-simplex $ABCD$.
\item
The ten edges of the $K_5$ $\{A,B,C,D,E\}$.
\end{enumerate}

If we can rule out 1 and 2 then 3 will follow, because
along with the $K_5$ we would have the five associated $3$-simplices:
Together these would form an $S^3$, not possible as a subcomplex
of the topological 3-ball $\tau$.

To rule out 1 and 2,
the strategy will be to consider a minimal counterexample,
and find a modification that would produce a smaller counterexample.

If $\sigma$ has a vertex of degree 3,
it is a connected sum with one summand a single tetrahedron.
We can pull off the vertex from $\sigma$, while removing the associated
tet from $\tau$, to produce a smaller counterexample.

Since $\sigma$ has no vertex of degree $3$, there are at least
$\maxdeg \geq 4$ disjointly eligible tets: tets that each meet the
boundary in two faces, with these pairs of faces disjoint.
Removing any one of these tets performs an `edge flip' on the boundary
$\sigma$,
which either (a) remains a triangulation of the 2-sphere,
or (b) becomes an almost disjoint union where
two such triangulations are joined along an edge.
We will show that some choice for the flip does not remove the taboo
configuration.
In case (a) we have a smaller counterexample already.
In case (b) the key is that if a $K_3$ or $K_4$ occurs
in the almost disjoint union, it occurs on one side or the other,
also giving a counterexample.

Configuration 1:
When we remove a tet from $\tau$ to flip an edge of $\sigma$,
the only edge that disappears from $\tau$ is the edge that gets flipped.
There are four choices of the edge to flip, and three edges we must
avoid flipping.

Configuration 2:
Now there are six edges we must avoid flipping. Having dealt with
configuration 1, we know that along with the six edges of $\{A,B,C,D\}$,
$\tau$ contains all four faces $ABC$, $ABD$, $ACD$, $BCD$. At most two of
these can be boundary faces, since $\sigma$ has no vertex of degree $3$;
so at most five of the six $K_4$ edges lie on the boundary. Treating the
octahedron separately as a special case, we may assume $\maxdeg \geq 5$,
which gives at least five boundary edges available to flip. The only
potential obstruction is when exactly five of the six $K_4$ edges lie on
the boundary, the interior edge being, say, $AD$; then the two boundary
faces of the $K_4$ are $ABC$ and $BCD$. Removing an eligible tet does no
harm unless one of its two boundary faces is $ABC$ or $BCD$. At most two
of our disjointly eligible tets can have $ABC$ or $BCD$ as a boundary
face, leaving at least $\maxdeg-2 \geq 3$ further eligible tets whose
removal flips an edge while preserving the $K_4$.
\proofend

\bibliography{taut}
\bibliographystyle{amsplain}
\end{document}
